\begin{statementbox}
\textbf{条件。} 在同一平面内，$A\ne B$ 为两定点，$\lambda>0$ 为常数。
动点 $P$ 满足 $PA/PB=\lambda$；比值要求 $P\ne B$，且此时也有 $P\ne A$。

% First-page model: a=1, lambda=2; all lengths use the same scale.
\begin{center}
\begin{tikzpicture}[font=\small]
\draw[black!30,rounded corners=2pt] (0,0) rectangle (14.4,8.8);
\node[anchor=west,font=\small\bfseries] at (.55,8.25)
  {先看图：到两定点的距离比为 $2$，动点走出一个圆周};
\node[anchor=west] at (.7,6.8) {$A(-1,0),\ B(1,0)$};
\node[anchor=west] at (.7,6.05) {$\color{blue}PA\color{black}=2\color{teal!80!black}PB$};
\begin{scope}[shift={(4.5,4.15)},scale=2.2]
\coordinate (A) at (-1.0000000000,0.0000000000);
\coordinate (B) at (1.0000000000,0.0000000000);
\coordinate (C) at (1.6666666667,0.0000000000);
\coordinate (D) at (0.3333333333,0.0000000000);
\coordinate (E) at (3.0000000000,0.0000000000);
\coordinate (P) at (1.6666666667,1.3333333333);
\draw[black!50] (-1.3,0)--(3.4,0);
\draw[blue,very thick] (C) circle (1.3333333333);
\draw[orange!85!black,very thick] (D)--(E);
\draw[blue,thick] (A)--(P);
\draw[teal!80!black,thick] (B)--(P);
\draw[orange!85!black,dashed,thick] (C)--(P)
  node[midway,right=5pt] {$R$};
\fill (A) circle (1.3pt) node[below=6pt] {$A$};
\fill (B) circle (1.3pt) node[below=15pt] {$B$};
\fill (C) circle (1.3pt) node[below=5pt] {$C$};
\fill (D) circle (1.3pt) node[below left=5pt] {$D$};
\fill (E) circle (1.3pt) node[below right=5pt] {$E$};
\fill (P) circle (1.5pt) node[above=5pt] {$P$};
\end{scope}
\node[anchor=west] at (.7,.63)
  {$D$ 内分、$E$ 外分：$DE$ 是直径，$C$ 是中点；$C(5/3,0),\ R=4/3$。};
\end{tikzpicture}
\end{center}

\textbf{先认图中的对象：} 蓝色圆周是轨迹，$C$ 是圆心，$DE$ 是直径；
$A,B$ 是给定的两点，\textbf{都不在圆周上}。

\textbf{标准坐标。} 取 $A(-a,0),B(a,0)$，$a>0$。
当 $\lambda\ne1$ 时，轨迹是\textbf{整个圆周}：
\[
 \boxed{(x-h)^2+y^2=R^2},\qquad
 h=\frac{a(\lambda^2+1)}{\lambda^2-1},\quad
 R=\frac{2a\lambda}{|\lambda^2-1|}>0.
\]
\textbf{分界：}$\lambda=1$ 时不是圆，轨迹为 $AB$ 的中垂线。
\end{statementbox}
\clearpage
\begin{statementbox}
\textbf{圆心位置与速算公式。}
圆心 $C(h,0)$ 在线段 $AB$ 外：$\lambda>1$ 在 $B$ 外侧，
$0<\lambda<1$ 在 $A$ 外侧。圆周不经过 $A,B$。
% Same metric in both axes; a=1 and lambda=2.
\begin{center}\begin{tikzpicture}[font=\small]
\draw[black!30,rounded corners=2pt] (0,0) rectangle (14.4,7.4);
\node[anchor=west,font=\small\bfseries] at (.55,6.85)
  {圆心与半径直接对照：$CA:CB=4:1$，$R^2=CA\cdot CB$};
\node[anchor=west] at (.75,4.9) {示例：$PA/PB=2$};
\node[anchor=west] at (.75,4.1) {圆心 $C$ 在 $B$ 外侧};
\node[anchor=west] at (.75,3.3) {$R=\sqrt{\frac83\cdot\frac23}=\frac43$};
\begin{scope}[shift={(6.3166666667,3.9)},scale=1.25]
\coordinate (A) at (-1.0000000000,0.0000000000);
\coordinate (B) at (1.0000000000,0.0000000000);
\coordinate (C) at (1.6666666667,0.0000000000);
\coordinate (D) at (0.3333333333,0.0000000000);
\coordinate (E) at (3.0000000000,0.0000000000);
\coordinate (P) at (1.6666666667,1.3333333333);
\draw[black!50] (-1.2,0)--(3.4,0);
\draw[blue,very thick] (C) circle (1.3333333333);
\draw[orange!85!black,thick] (D)--(E);
\draw[orange!85!black,dashed,thick] (C)--(P)
  node[midway,right=5pt] {$R$};
\fill (A) circle (1.3pt) node[below=6pt] {$A$};
\fill (B) circle (1.3pt) node[below=15pt] {$B$};
\fill (C) circle (1.3pt) node[below=5pt] {$C$};
\fill (D) circle (1.3pt) node[below left=5pt] {$D$};
\fill (E) circle (1.3pt) node[below right=5pt] {$E$};
\fill (P) circle (1.3pt) node[above=5pt] {$P$};
\end{scope}
\draw[blue,<->,thick] (5.0666666667,.45)--(8.4,.45)
  node[midway,above=3pt] {$CA=8/3$};
\draw[teal!80!black,<->,thick] (7.5666666667,1.25)--(8.4,1.25)
  node[midway,above=3pt] {$CB=2/3$};
\draw[black!25,densely dotted] (5.0666666667,3.25)--(5.0666666667,.6);
\draw[black!25,densely dotted] (7.5666666667,2.8)--(7.5666666667,1.9);
\draw[black!25,densely dotted] (8.4,2.0)--(8.4,1.9);
\draw[black!25,densely dotted] (8.4,1.05)--(8.4,.6);
\end{tikzpicture}\end{center}


\textbf{最快求圆心、半径：平方外分，距离开方。}
一般地，先按\textbf{平方比外分}找 $C$，再由两段长度求半径：
\[
 \boxed{CA:CB=\lambda^2:1}\qquad\text{（平方外分）},
 \qquad \boxed{R=\sqrt{CA\cdot CB}}\qquad\text{（距离开方）}.
\]
这两式中 $CA,CB$ 均表示\textbf{圆心到定点}的长度，并非 $PA,PB$。

\textbf{几何等价表述。} 直线 $AB$ 上满足同一距离比的内分点 $D$、外分点 $E$ 为
\[
 D\left(\frac{a(\lambda-1)}{\lambda+1},0\right),\qquad
 E\left(\frac{a(\lambda+1)}{\lambda-1},0\right).
\]
轨迹恰为\textbf{以 $DE$ 为直径的圆}，故 $C$ 为 $DE$ 中点，$R=DE/2$。

\textbf{一般坐标。} 若 $A(x_1,y_1),B(x_2,y_2)$，记 $L=AB>0$，则
\[
 C\left(\frac{x_1-\lambda^2x_2}{1-\lambda^2},
         \frac{y_1-\lambda^2y_2}{1-\lambda^2}\right),\qquad
 R=\frac{\lambda L}{|1-\lambda^2|}.
\]
这些圆公式均要求 $\lambda>0,\lambda\ne1$。
\textbf{当 $\lambda=1$ 时}，轨迹为 $AB$ 的中垂线；在标准坐标下为 $x=0$。
本条讨论平面内的自由动点；若另加直线等限制，应取圆与限制集合的交集。
\end{statementbox}
